\pagenumbering{gobble}

\begin{titlepage}
\centering
{\large \textbf{Universidad:} \textcolor{red}{pendiente completar nombre institucional oficial} \par}
\vspace{0.6cm}
{\large \textbf{Programa doctoral:} \textcolor{red}{pendiente completar denominacion formal} \par}
\vspace{0.3cm}
{\large \textbf{Version del documento:} borrador de revision interna doctoral \par}
\vspace{1.6cm}
{\Huge \textbf{\ThesisTitle} \par}
\vspace{0.5cm}
{\Large Monografia doctoral de investigacion \par}
\vspace{0.4cm}
{\large Tesis sobre optimizacion de memoria y ejecucion heterogenea en entrenamiento profundo \par}
\vspace{1.2cm}
{\large Doctorando: \ThesisAuthor \par}
\vspace{0.4cm}
{\large Direccion: \textcolor{red}{pendiente completar director(a) y codirector(a)} \par}
\vspace{0.4cm}
{\large Lugar y fecha de presentacion: \textcolor{red}{pendiente completar ciudad, pais y mes de deposito} \par}
\vfill
{\large Fecha de compilacion: \today \par}
\end{titlepage}

\clearpage
\pagenumbering{roman}

\chapter*{Resumen}
\addcontentsline{toc}{chapter}{Resumen}
Esta monografia estudia el problema de la particion heterogenea CPU-GPU para entrenamiento profundo como un problema metodologico y experimental completo, no solo como una heuristica de asignacion. La tesis sostiene que una politica de particion solo adquiere valor cientifico cuando puede fundarse en profiling empirico reproducible, traducirse a una formulacion declarativa con restricciones fisicas explicitas y cerrarse en simulacion y ejecucion observada sobre runtime real. Sobre esa base, el manuscrito construye un pipeline en cinco fases que enlaza medicion por capa, agregacion robusta, optimizacion ILP, representacion formal del plan y validacion hibrida con trazabilidad de artefactos.

Los resultados disponibles muestran que el proyecto ya dispone de cierre metodologico funcional: mide, decide, simula y ejecuta. El profiling produce coeficientes auditables de tiempo, energia, memoria y transferencia; el ILP encuentra politicas superiores a baselines simples en casos controlados y campañas smoke; y el runtime hibrido permite contrastar el comportamiento previsto frente al observado, incluyendo evidencia focal de extensiones avanzadas como separacion forward/backward y prefetching. La principal brecha restante no es de implementacion sino de amplitud empirica: falta consolidar la campaña multi-servidor y ampliar la matriz final de validacion para convertir la evidencia focal actual en cierre doctoral transversal.

\vspace{0.6cm}
\noindent\textbf{Palabras clave:} particion heterogenea CPU-GPU; entrenamiento profundo; optimizacion ILP; profiling empirico; gestion de memoria; simulacion y runtime hibrido.

\chapter*{Abstract}
\addcontentsline{toc}{chapter}{Abstract}
This dissertation studies heterogeneous CPU-GPU partitioning for deep learning training as a full methodological and experimental problem rather than as an isolated placement heuristic. The central claim is that a partitioning policy only becomes scientifically meaningful when it is grounded on reproducible empirical profiling, translated into a declarative optimization model with explicit physical constraints, and finally validated through simulation and observed hybrid execution on a real runtime. Accordingly, the manuscript develops a five-phase pipeline that links per-layer measurement, robust aggregation, ILP optimization, formal plan representation, and runtime validation under a traceable evidence contract.

Current results already demonstrate methodological closure: the project can measure, decide, simulate, and execute. Profiling yields auditable coefficients for time, energy, memory, and transfer; the ILP improves upon simple baselines in smoke and controlled settings; and the hybrid runtime supports prediction-versus-observation analysis, including focused evidence for advanced features such as forward/backward separation and explicit prefetching. The main remaining gap is therefore empirical breadth rather than software capability: a final multi-server campaign is still required to transform the current focal evidence into a fully doctoral-grade cross-platform validation.

\vspace{0.6cm}
\noindent\textbf{Keywords:} heterogeneous CPU-GPU partitioning; deep learning training; integer linear programming; empirical profiling; memory management; hybrid runtime validation.

\clearpage

\chapter*{Convenciones del manuscrito}
\addcontentsline{toc}{chapter}{Convenciones del manuscrito}
El documento distingue con rigor entre capacidad implementada, validacion focal y cierre empirico amplio. Cuando una afirmacion describe una capacidad del sistema ya materializada en el repositorio, se formula en presente. Cuando una afirmacion depende de evidencia controlada disponible pero todavia no generalizada a una matriz experimental completa, se explicita su alcance focal. Cuando una conclusion requiere una campaña final aun no ejecutada, se marca como trabajo de cierre pendiente y no como resultado ya demostrado. Esta convencion evita inflar el alcance de la evidencia y mantiene la coherencia epistemologica del manuscrito.

\chapter*{Alcance de esta version}
\addcontentsline{toc}{chapter}{Alcance de esta version}
La presente version debe leerse como un borrador doctoral avanzado de preentrega interna. Su estructura argumental, su cadena metodologica y su soporte empirico principal ya estan integrados, pero todavia quedan por completar ciertos elementos propios de la version de deposito: datos institucionales definitivos de portada, consolidacion de la campana multi-servidor, ampliacion bibliografica comparada y ajuste final de algunos cuadros y figuras. Esta explicacion no debilita el manuscrito; delimita con precision el alcance de lo ya demostrado y preserva la honestidad epistemologica del documento. En particular, los campos institucionales marcados en rojo no responden a una carencia de estructura del manuscrito, sino a la ausencia de informacion oficial cerrada dentro del repositorio de trabajo.

\chapter*{Siglas y abreviaturas}
\addcontentsline{toc}{chapter}{Siglas y abreviaturas}
\small
\begin{tabularx}{\textwidth}{lX}
CPU & Unidad Central de Procesamiento. \\
GPU & Unidad de Procesamiento Grafico. \\
ILP & Programacion Lineal Entera. \\
VRAM & Memoria de video residente en el acelerador. \\
DAG & Grafo aciclico dirigido. \\
DDP & Paralelismo de datos distribuido. \\
FSDP & Paralelizacion completamente fragmentada de parametros y estados. \\
\end{tabularx}
\normalsize